% !PRINTED 605
\weberpage{605}

In this, $\chi$ may be any of the characters $\chi_\kappa$, and the sum extends over all ideals $\mathfrak{a}$ of $\overline{O}$. By (11), if we set

\[
G_\kappa = \sum^{i}\chi_\kappa(\overline{A}_i)\,C_i
\]

then

\weq{17}{\begin{aligned}
Q_1 &= \frac{hG}{s-1} + G_1, \\
Q_\kappa &= G_\kappa, \quad \kappa > 1.
\end{aligned}}
% [source sic: equation (17) prints hG, whereas equation (21) later uses hg for the same residue]

where the $G_1, G_2, \dots, G_h$ are \gesperrt{functions of $s$ which retain a finite value for $s = 1$.}

The sums $Q$ can be transformed in the same way that we have already come to know in \S\,197 of the second volume.

Accordingly, if $\mathfrak{p}$ denotes any prime ideal in $\overline{O}$:

\[
\begin{aligned}
\frac{1}{1 - \chi(\mathfrak{p})N(\mathfrak{p})^{-s}} &= 1 + \chi(\mathfrak{p})N(\mathfrak{p})^{-s} + \chi(\mathfrak{p})^2 N(\mathfrak{p})^{-2s} + \cdots \\
&= 1 + \frac{\chi(\mathfrak{p})}{N(\mathfrak{p})^s} + \frac{\chi(\mathfrak{p}^2)}{N(\mathfrak{p}^2)^s} + \frac{\chi(\mathfrak{p}^3)}{N(\mathfrak{p}^3)^s} + \cdots,
\end{aligned}
\]

and since all ideals $\mathfrak{a}$ can be represented as products of such ideals $\mathfrak{p}$, it follows that:

\weq{18}{Q = \Pi\,\frac{1}{1 - \chi(\mathfrak{p})N(\mathfrak{p})^{-s}},}

where the infinite product $\Pi$ extends over all prime ideals $\mathfrak{p}$ in $\overline{O}$.

To every ideal $\mathfrak{a}$ in $\overline{O}$ there belongs a certain exponent $\varphi$, that is, there is a certain lowest positive exponent $\varphi$ for which $\mathfrak{a}^\varphi$ is contained in the principal class $\overline{A}_1$. The number $\varphi$ is always a divisor of the class number $h$. We set:

\weq{19}{h = \varepsilon\varphi.}

If in this sense $\mathfrak{a}$ belongs to the exponent $\varphi$, then all the characters $\chi_k(\mathfrak{a})$ are $\varphi$th roots of unity, and every $\varphi$th root of unity occurs $\varepsilon$ times among them (Bd.\ II, \S\,198, 1.).
% [source sic: the index switches here from \chi_\kappa to \chi_k]

If, then, $\mathfrak{p}$ belongs to the exponent $\varphi$, we have

\[
\begin{aligned}
&[1 - \chi_1(\mathfrak{p})N(\mathfrak{p})^{-s}][1 - \chi_2(\mathfrak{p})N(\mathfrak{p})^{-s}]\cdots[1 - \chi_h(\mathfrak{p})N(\mathfrak{p})^{-s}] \\
&\qquad = [1 - N(\mathfrak{p})^{-\varphi s}]^\varepsilon,
\end{aligned}
\]

and from (18) it follows that:

\weq{20}{Q_1 Q_2 \cdots Q_h = \Pi\,\frac{1}{[1 - N(\mathfrak{p})^{-\varphi s}]^\varepsilon}.}
